\originalchapter{麦克斯韦方程与电磁能流}\label{ch:maxwell}
\originalsection{位移电流与统一方程}
充电电容器的导线有电流, 极板间真空却无真实电荷穿过. 对同一安培环路, 选横穿导线的面得到$I$, 改选从极板间鼓出的面得到零, 旧稳恒式会矛盾. 极板间电通量随时间增长, 且对整个板面积$\eps\dot\Phi_E=\dot Q=I$. 麦克斯韦补充位移电流$I_d=\eps\dot\Phi_E$, 使变化电场也能贡献磁场环量.

在真空, 以总电荷密度和总电流密度为源, 四个积分方程为
\begin{align}
\oint_S\ee\cdot\dd\vct A&=Q_{\rm enc}/\eps,&\oint_S\bb\cdot\dd\vct A&=0,\nonumber\\
\oint_C\ee\cdot\dd\vct\ell&=-\frac{\dd}{\dd t}\int_{S_C}\bb\cdot\dd\vct A,&
\oint_C\bb\cdot\dd\vct\ell&=\mu_0I_{\rm enc}+\mu_0\eps\frac{\dd}{\dd t}\int_{S_C}\ee\cdot\dd\vct A.\label{eq:maxwell}
\end{align}
前两式用闭合面$S$, 后两式用固定闭合曲线$C$及以它为边界的有向面$S_C$. 第一式描述电荷源, 第二式描述无磁单极源, 第三式描述变化磁场激起电场, 第四式描述传导电流与变化电场激起磁场. 稳恒时第四式退回安培定律, 第三式退回静电场环量零.

在场足够光滑处, 利用散度与旋度的积分关系可写
\[
\nabla\cdot\ee=\rho/\eps,\quad\nabla\cdot\bb=0,\quad
\nabla\times\ee=-\partial_t\bb,\quad\nabla\times\bb=\mu_0\jj+\mu_0\eps\partial_t\ee.
\]
对最后一式取散度, 左边恒为零, 结合第一式得到$\nabla\cdot\jj+\partial_t\rho=0$, 正好是电荷守恒. 位移电流不是真空中新的粒子流, 而是使场规律与守恒相容的变化电场项.

\originalsection{平面波与极化}
波动方程也可从已经会用的积分环路求得, 不必先做三维旋度运算. 考虑场只随$x,t$变化且$\ee=E_y\uv y$, $\bb=B_z\uv z$. 在$xy$平面取小矩形, 正法向$+z$, 逆时针绕行, 两条沿$y$边给环量$[E_y(x+\Delta x)-E_y(x)]\Delta y$. 法拉第式除以面积, 缩小矩形得
\[
\partial_xE_y=-\partial_tB_z.
\]
在$xz$平面取正法向$+y$的小矩形, 相应正绕行使沿$z$的环量为$-[B_z(x+\Delta x)-B_z(x)]\Delta z$. 安培--麦克斯韦式给
\[
-\partial_xB_z=\mu_0\eps\partial_tE_y.
\]
对第一式再对$x$求导, 将第二式对$t$求导后代入, 得$\partial_x^2E_y=\mu_0\eps\partial_t^2E_y$. 这里需要场二阶偏导连续, 以便交换求导次序.

试代任意平移波形$E_y=F(x-vt)$, 链式法则给$\partial_x^2E_y=F''$、$\partial_t^2E_y=v^2F''$, 对有非零曲率的波形得到$v^2=1/(\mu_0\eps)$. 固定波形某一点$x-vt=\xi_0$随时间以$x=vt+\xi_0$向前移, 所以$F(x-ct)$沿正$x$传播, $F(x+ct)$沿负$x$传播. 周期波固定相位$kx-\omega t$也给出速度$\omega/k=c$.

平面波若另有$E_x$, 无源高斯定律要求$\partial_xE_x=0$; 安培--麦克斯韦式的$x$分量又给$\partial_tE_x=0$. 因而它只能是另加的静态常量背景, 不属于传播的振荡场. 磁场同理, 所以无背景的平面波为横波.

无源真空中对法拉第式取旋度, 用$\nabla\times(\nabla\times\ee)=\nabla(\nabla\cdot\ee)-\nabla^2\ee$和安培--麦克斯韦式, 得$\nabla^2\ee=\mu_0\eps\partial_t^2\ee$, 磁场同样满足波动方程. 因而波速$c=1/\sqrt{\mu_0\eps}$, 与光速相同. 这把光和电磁作用联系起来, 不表示所有材料中的传播速度均为$c$.

沿正$x$方向传播的一组解为
\[
\ee=E_0\sin(kx-\omega t)\uv y,\qquad
\bb=\frac{E_0}{c}\sin(kx-\omega t)\uv z,\qquad \omega=ck.
\]
代入$\partial_xE_y=-\partial_tB_z$可检查振幅与符号. 场彼此垂直, 都垂直传播方向, 是横波. 波长$\lambda=2\pi/k$, 频率$f=\omega/(2\pi)$, 满足$c=f\lambda$. 任意传播单位矢量$\uv p$的真空平面行波有$\bb=(\uv p\times\ee)/c$, $\ee\times\bb$沿传播方向.

线偏振方向是电场振动的轴, 场每半周期反向不改变偏振轴. 两个相互垂直且有固定相位差的分量也可合成圆或椭圆偏振. 理想线偏振片只保留沿透射轴的电场分量, 所以线偏振入射光通过转角$\alpha$后的强度为$I_0\cos^2\alpha$. 无偏振入射光在各方位平均给出$I_0/2$. 水平表面的反射眩光常有较强水平偏振分量; “所有散射都只保留一个方向”并非普遍结论.

\originalsection{能量守恒与坡印廷矢量}
光照能加热物体, 说明传播的场携带能量. 场能密度和能流密度为
\[
u=\frac{\eps E^2}2+\frac{B^2}{2\mu_0},\qquad
\vct S=\frac1{\mu_0}\ee\times\bb.
\]
两者单位分别为J/m$^3$与W/m$^2$. 用恒等式$\nabla\cdot(\ee\times\bb)=\bb\cdot(\nabla\times\ee)-\ee\cdot(\nabla\times\bb)$代入场方程, 得
\begin{equation}\label{eq:poynting}
\partial_tu+\nabla\cdot\vct S=-\jj\cdot\ee.
\end{equation}
积分后表示体内场能增加率、向外能流与场对物质做功之间的平衡. 电阻中$\jj\cdot\ee>0$, 电池输出能量的区域可能为负. 这比“能量只在导线里面走”更完整.

行进平面波中两种能量密度相等, 且$S=cu$. 正弦波平均强度$I_{\rm light}=\langle S\rangle=\eps cE_0^2/2$. 这里$I_{\rm light}$是光强, 单位W/m$^2$, 不同于电路电流. 对正入射理想吸收面光压为$I_{\rm light}/c$, 对正入射理想反射面为$2I_{\rm light}/c$, 原因是每单位能量的电磁动量为$1/c$, 反射时动量方向反转.

\originalsection{元件中的能流与波的产生}
在圆柱电阻中, 电场沿电流, 磁场绕轴, $\ee\times\bb$指向轴, 能量通过侧面进入电阻并变成热. 半径$a$, 长$\ell$, 场$E=V/\ell$, 表面$B=\mu_0I/(2\pi a)$, 所以向内功率为$EB(2\pi a\ell)/\mu_0=VI$. 充电电容中位移电流产生环向磁场. 半径$a$圆板间距$d\ll a$, 变化时间尺度$\tau\gg a/c$, 可在忽略边缘的慢变模型中用均匀有符号场$E(t)\uv z$. 传入正极板的电流$I=\eps\pi a^2\dot E$. 半径$s<a$安培圆盘的位移电流为$\eps\pi s^2\dot E$, 所以$B_\phi=\mu_0\eps s\dot E/2$, $S_s=-\eps sE\dot E/2$. 两端面能流法向分量零, 侧面向内功率为
\[
-S_s(a)2\pi ad=\eps\pi a^2dE\dot E
=\frac{\dd}{\dd t}\left(\frac{\eps E^2}2\pi a^2d\right).
\]
它正是场能增加率. 判断流入还是流出应看$E\dot E$, 不是只看$\dot E$: 同方向场强增大时流入, 储能下降时流出. 反号电压的充电也可以使$E<0,\dot E<0$, 两者乘积仍正. 第22章按原题逐问核同一模型. 有限几何的边缘和外部场不能一概删去, 本算例只在上述近似范围内使用.

长螺线管电流增大时, 内部$B=\mu_0nI$增长. 半径$a$内, 法拉第定律给出$E_\phi2\pi s=-\pi s^2\dot B$, 即$E_\phi=-s\dot B/2$. 因而$S_s=E_\phi B/\mu_0=-sB\dot B/(2\mu_0)$. 当$B\dot B>0$时向内, 侧面输入功率$\pi a^2\ell B\dot B/\mu_0=\dd[\pi a^2\ell B^2/(2\mu_0)]/\dd t$. 衰减时能流反向. 该理想模型要求变化时间远大于相关长度除以光速, 忽略端部、辐射及较小的电场储能贡献, 因而属于慢变近似; 在此范围直接核对了磁能的来去, 不把它当成任意快速变化电流的精确解.

加速电荷可向外产生电磁波. 对位于$x=0$的无限均匀面, 两半空间均真空, 没有入射波, 设沿$y$有规定的面电流$K(t)$. 仅取向两侧外发射的横向场. 向两侧发出的行波在面处有相同切向电场$E_y$, 右侧$B_z=E_y/c$, 左侧$B_z=-E_y/c$. 跨面的磁场边界式给出$B_z(0^+)-B_z(0^-)=-\mu_0K$, 所以$E_y=-Z_0K/2$, $Z_0=\mu_0c$. 两侧总辐射功率每单位面积为$2E_y^2/Z_0=Z_0K^2/2$, 正好等于维持电流抵抗反作用场所需功率$-KE_y$. 这个边界负号可从$\uv x\times(\bb_+-\bb_-)=\mu_0K\uv y$验证: $\uv x\times\uv z=-\uv y$. 面外场用推迟时间表达为
\[
E_y(x,t)=-\frac{Z_0}2K(t-|x|/c),\qquad
B_z(x,t)=\operatorname{sgn}(x)E_y(x,t)/c.
\]
所以远处不与电流源同时响应. 若$K=K_0\sin\omega t$, 两侧周期平均总辐射功率每面积是$Z_0K_0^2/4$, 与瞬时表达的系数不同. 用“摇动带电面”的图像时,$K=\sigma\dot y$, 还叠加静态法向场$E_x^\pm=\pm\sigma/(2\eps)$. 上述计算针对横向辐射和向左右平面的法向功率, 没有把静态库仑场误称行波. 源维持往复运动所做的功与辐射功率相符.

有限天线中振荡电流使电荷分离周期改变, 向外辐射; 靠近天线的场和远处行波不能混同. 常见谐振偶极天线两臂各约$\lambda/4$, 单臂与导电地面构成的单极天线长约$\lambda/4$. 长度关系来自电流、电压驻波的边界条件及传播相位, 不是金属电子漂移接近光速. 本册解释波的产生与能流, 远区辐射的完整多极推导见电动力学册.


简化天线的驻波边界可解释长度. 开路端没有电流流出, 是电流节点; 中央馈电附近近似电流腹. 正弦驻波从腹到相邻节点相位差为$\pi/2$, 距离为$\lambda/4$, 因而两臂总长约$\lambda/2$. 真实长度还受导体端部效应、附近材料与有效传播速度影响. 火花隙发射器先充电, 击穿后电容和线段电感产生衰减振荡, 能量逐步转为辐射与电阻损耗; 这连接第\ref{ch:ac}章的自由RLC与波的产生.

沿天线轴的往复电流使远区辐射偏振与该轴的投影有关. 短电偶极辐射在轴向很弱, 在垂直轴的方向较强, 其近似强度角因子为$\sin^2\theta$; 该短偶极式不等于任意半波天线的精确方向图. 轴向场的变化不能直接用某一平面波图判断全空间辐射. 接收线状天线沿电场方向时耦合较强, 转成垂直时响应减弱, 因而实验可同时辨别偏振与方向强弱.
